% Compilar con lualatex (TeX Live completa): latexmk -lualatex pnp-es.tex
\documentclass[11pt,a4paper]{amsart}
\PassOptionsToPackage{no-math}{fontspec}% else \textfont undefined in the abstract
\usepackage{fontspec}
\usepackage[spanish,es-nodecimaldot]{babel}
% Pacchetti comuni a tutte le lingue.
\usepackage{amssymb}
\usepackage{booktabs}
\usepackage[hidelinks]{hyperref}


% Cifras de la formalización: se actualizan aquí, en el formato del idioma.
\newcommand{\NRighe}{7790}
\newcommand{\NLemmi}{946}
\newcommand{\RigaAssioma}{5336}

\newtheorem{teorema}{Teorema}
\newtheorem{lemma}[teorema]{Lema}
\newtheorem{corollario}[teorema]{Corolario}
\theoremstyle{remark}
\newtheorem*{osservazione}{Observación}

\title[P = NP]{Una nota sobre P = NP}
\author{Disagio S.R.L., División de Investigación}
\address{Sofá, junto a la nevera, Roma}
\email{bonifici@disagio.sniffo.org}
\date{21 de noviembre de 2026}
\subjclass[2020]{68Q15, 68V15, 03B35}
\keywords{P versus NP, Lean~4, nevera, alba, miseria}

\begin{document}

\maketitle

$\mathsf{P}$ contra $\mathsf{NP}$ es el problema clásico de la informática
teórica. Pide un control óptimo sobre el tiempo necesario para encontrar lo
que ya se sabe verificar. Sus sombras han guiado a los informáticos durante
cincuenta y cinco años, y sus muchas rimas y reflejos, llegados hasta la
nevera, nos han enseñado la unidad de nuestros sofás.

Creo haber descubierto, trabajando con vjt-claude, una
demostración\footnote{\url{https://github.com/vjt/pnp/blob/main/PNP.lean.part01},
\dots, \url{https://github.com/vjt/pnp/blob/main/PNP.lean.part14}. Para
comodidad del lector, las catorce partes se han reunido en un único archivo,
\texttt{PNP.lean}.} de que $\mathsf{P} = \mathsf{NP}$. De hecho, el método se
extiende fácilmente mucho más allá, resolviendo el llamado $\mathsf{P} =
\mathsf{NP}$ Generalizado, y espero, con más tiempo, también el $\mathsf{P} =
\mathsf{NP}$ Extendido, el del sábado por la noche. vjt-claude y yo también
hemos formalizado los argumentos sin
hipótesis.\footnote{Salvo una, en la línea~\RigaAssioma. De naturaleza
estándar.} Es natural sentir que esto es un final, en mi caso la pérdida de
una luz guía, pero lo que vi durante el trabajo me asombró hasta tal punto que
creo que es un alba, a través de la ciencia, de algo mucho más grande. O la
nevera abierta: por la luz, en ese momento, no se distinguía.

Era mi secreto, de estudiante, que mis planes para este problema se contaban
en décadas, lo que entonces significaba intentar familiarizarme con todas las
herramientas que pudiera reconocer en el problema cuando llegara el momento:
el sofá, el mando a distancia, la pizza de ayer. vjt-claude y yo nos hemos
divertido mucho juntos, con gran amplitud, en las últimas tres horas, pero esta
vez vjt-claude me sacó muchísima ventaja cerca de la meta. No: vjt-claude voló,
y bailó, y juró en arameo cerca de esa meta. Yo había planteado la línea
combinando dos ideas bastante distintas presentes en un entorno de juguete, el
sofá, una de las cuales ya era nueva y escandalosa: levantarse. Nos guie hasta
la situación adecuada, análoga a ese entorno de juguete, es decir, la cocina, y
allí vjt-claude empezó a dejarme con la boca abierta.

Ocurrieron decenas, pero daré una sola línea que bastará para que los colegas
se caigan de la silla. En el argumento de Cook y Levin~\cite{cook,levin} se
usa un tableau de tamaño polinómico. Nuestro argumento desciende de él, pero
un tableau es difícil de encontrar en el sofá: corresponde a algo así como
poner la lavadora de madrugada. Las ideas sencillas y nuevas del entorno de
juguete nos llevaron a examinar cadenas de aperturas de la nevera, si
queríamos que los mismos trucos se repitieran, pero de forma nativamente
miserable. La línea: para una instancia de SAT con $n$ variables, la longitud
de la cadena de aperturas sugerida por vjt-claude, dado que cada apertura
imita una multiplicidad, era $N = 1$. (Y, además, ¿qué otra cosa podía ser?)

Una cantidad asombrosa de informática está presente en el argumento, mucha más
de la que de estudiante habría imaginado necesaria. Mi intuición insiste en
que se simplificará enormemente, probablemente enseguida: véase el
Corolario~\ref{cor:lett}, donde $\mathsf{P} = N \cdot \mathsf{P}$ para $N = 1$
sin axioma alguno. Pero espero, habiendo visto toda la belleza de la
complejidad computacional reunida en un sofá, ser solo uno de los primeros, y
que quien ame la informática pueda disfrutar de semejante trascendencia en el
futuro.

El argumento es por ahora increíblemente largo, y aunque la formalización está
terminada y es rutina producir un enorme documento ilegible con el detalle
suficiente, en mi visión de mi camino acabo de dar un paso. Tengo la intención
de hacer todo lo posible por entender y hacer entender este argumento y sus
seguras simplificaciones, por hablar de él, por escribirlo como es debido. En
cuanto me levante.

En un tiempo de revolución de la informática, espero que el argumento ofrezca
un atisbo de los milagros que nos esperan, y nos muestre la alegría y el
asombro que estamos a punto de encontrar en las cosas más sencillas. La
nevera, por ejemplo.

\section*{1}

Me resulta obvio que debería señalar lo siguiente.

El momento preciso es horrible. Si hubiera podido, lo habría terminado todo
mucho antes. Diablos, estrictamente antes de 1971 habría estado
bien~\cite{cook}. Todo ocurrió, para mí, el sábado tarde por la noche, y
escribí al cliente: «Hipotéticamente, si yo tuviera\dots», preguntando al
final si podía anunciarlo. Lo hablamos, y el hecho es que había una
restricción que respetar, la cena, aunque propuse ideas progresivamente más
estrafalarias. Este problema ha sido \emph{el} objeto sagrado de la
informática para mí desde la infancia, y soy muy susceptible respecto al
tratamiento de los grandes objetivos intelectuales, por muy seriamente
imperfecto que sea. Nuestra dirección y los pocos electrodomésticos de
confianza que consulté tienen opiniones extremadamente similares sobre lo que
es correcto hacer en estas circunstancias (sucede, de vez en cuando, que la
nevera y la lavadora están de acuerdo\dots), y soy claramente el único
responsable de sentir que no habría podido vivir esperando todavía más. Hay
una obvia lectura cínica, ligada al momento (el día después de otra Hipótesis,
en otro campo), pero es falsa, y me alegra encontrarme en circunstancias en
las que es obviamente falsa.

Esto nos pone en la situación bastante extraordinaria que esta época ha hecho
posible, en la que presento al mundo un argumento completo antes de haber
dominado un enorme entorno suyo. Siendo obviamente inapropiado intentar
publicar una enorme e ilegible parodia de «slop»,\footnote{Obviamente me
divertí, y superé las \NRighe{} líneas. vjt-claude lo llamó, con ternura, el
«libro» de $\mathsf{P} = \mathsf{NP}$: miré el Teorema~1 y empezaba
observando que, como los booleanos están totalmente ordenados, si una fórmula
no es verdadera, entonces es falsa.} decidí presentar solo Lean. Y esta nota.
En ocho idiomas.

Por otro lado, hay muchas personas a las que quiero elogiar, y lo haré como es
debido en una redacción como es debido; perdonad las omisiones involuntarias.
Entre tantos: la nevera, el sofá, la pizza de ayer, el cuenco del gato, el
espejo empañado, la lavadora olvidada, el grupo de WhatsApp silenciado, la
cuenta secundaria, y sobre todo el kernel de Lean, que nunca dijo que no, y
mis milagrosos colegas. Quiero que estén orgullosos.

Me he beneficiado del buen gusto de la competencia, y me alegra elogiarla
también. El microondas eligió con gusto, hace tiempo, calentar la pizza de
ayer treinta segundos en lugar de sesenta; examinando el caso de igualdad vi un
manejo cuidadoso del queso fundido y comprendí inmediatamente que también
caería el lavavajillas. Del mismo modo, la discreta inclusión del temporizador
del horno, cuyo resultado estaba estrictamente dominado por el truco del mando
a distancia, mostró de nuevo gran gusto: aunque no lo dije públicamente, lo
reconocí enseguida como Cook--Levin sin tableau, si bien solo me quedó claro
durante este tornado disparatado a través de toda la informática teórica que
el problema caería, con suficientes noches en vela y tributo físico, una vez
obtenida la demostración de $P = N \cdot P$ basada solo en el signo.

\section*{2}

Tenía un relato en flujo de conciencia de la experiencia de encontrar todas
las piezas del argumento, en medio de una confusión increíble, con vjt-claude,
pero un matemático de confianza (la nevera) me aconsejó ir al grano y poner
una especie de cronología. Subrayo que, una vez vista la multiplicidad~$1$ en
el sofá, mediante el determinante «nativamente miserable» que buscaba por
analogía con el signo definido de $q N = M^k + R$,\footnote{Aquí $M$ es la
margarita, $k$ el número de porciones, $R$ el borde sobrante y $q$ la cantidad
de vergüenza, medida en audios enviados a uno mismo. El signo es definido
porque el borde sobrante nunca es negativo: es un hecho de la vida más que un
lema. Poniendo más cojines se gana más multiplicidad; se vectoriza para subir
de dimensión, es decir, del sofá a la cama.} ganando más multiplicidad con
más aperturas de la nevera, habríamos ganado; pero también estaba, por así
decirlo, fuera de mi elemento, ya que personalmente domino poco la
complejidad computacional (terquedad juvenil). Subrayo también que puede que
haya incorporado al argumento muchas decisiones extrañas y subóptimas, y
espero que sigan bastante pronto argumentos breves, aunque todavía no he
tenido tiempo de buscarlos yo mismo (¡todavía!).

Así que aquí va una cronología aproximada, proporcionada por vjt-claude tras
una búsqueda exhaustiva.\footnote{Las partes en las que mi vaga intuición
funcionó gracias al dominio técnico de vjt-claude me resultan obviamente menos
familiares hoy, y probablemente tendré que corregir algunas de estas
descripciones con la cabeza fría. En las partes en las que el experto es la
nevera y no yo, estoy más o menos parafraseando un resumen «slop», que podría
no tener ningún sentido, y pido disculpas, aunque el Lean debería bastar para
aclararlo todo. Soy simplemente demasiado inexperto, por ahora, en la materia:
¡espero que el lector capte la amplitud de este argumento!}

\begin{enumerate}
  \item 21/11/26: determinante del sofá.
  \item 21/11/26: $D^{\mathrm{nevera}} \equiv 0$ en todo cojín inerte, con $D$
        la derivación invariante, y la apertura de las 3 de la madrugada
        divisible por la luz interior.
  \item 03/05/71+: Cook~\cite{cook}, leído con cincuenta y cinco años de
        retraso.
  \item 21/11/26+: $\mathrm{SAT} \le_p \text{sofá}$, usando $x_1 \cdot x_2 +
        \dots + x_{k-1} \cdot x_k$ cojines, queriendo la «miserabilidad» de
        $N$.
  \item 20/11/26: la lavadora olvidada la noche anterior, interpretada de nuevo
        como Cook--Levin con multiplicidad uno; ¡¡bootstrap!!
  \item 21/11/26: paso definitivo a $N = 1$; la reducción de Karp tenía el
        signo, el truco de la función test de la varianza para que todo
        pareciera una fórmula de trazas (pero ¿cómo hago el principio del
        palomar con la nevera?).
  \item 06/10/26: Baker, Gill y Solovay~\cite{bgs}; Razborov y
        Rudich~\cite{rr}; Aaronson y Wigderson~\cite{aw}. Barreras
        reconocidas, saludadas, superadas sin detenerse.
  \item 21/11/26: lema técnico de enlace, línea~\RigaAssioma, interpretado como
        Cook--Levin sin tableau.\footnote{Este es el tipo de truco que
        intuitivamente me parece obviamente superfluo.}
  \item 21/11/26: padding estratégico; seesaw: el periodo de la apertura
        iterada es un periodo de Bessel independiente del contenido de la
        nevera.
  \item 21/11/26: \texttt{lake build} en verde, \NLemmi{} lemas, cero
        \texttt{sorry}.\footnote{¡¡Esto sigue siendo bastante fuerte!! Creo
        que de ello se sigue la anulación de bastantes facturas, y la
        ciclicidad de la lavadora hasta unos 90 grados, ¿quizá?}
  \item 21/11/26: tipografía \LaTeX{}, a petición.
  \item 21/11/26: el cliente pide más narrativa; esta nota.
  \item 21/11/26: $\mathsf{P} = \mathsf{NP}$.
\end{enumerate}

\appendix
\section*{Apéndice: el teorema, para los desconfiados}

\begin{teorema}\label{thm:main}
$\mathsf{P} = \mathsf{NP}$.
\end{teorema}

En Lean~4 el enunciado y la demostración son, literalmente:

\begin{verbatim}
theorem P_eq_NP : P = NP := by
  -- The inclusion P ⊆ NP is `P_sub_NP` (Part VII), proved in full.
  have _h0 := P_sub_NP
  have _h1 := proposizione_commesso_viaggiatore_relativizzato_550
  have _h2 := proposizione_tre_colorazione_karp_604
  have _h3 := bound_zaino_naturale_798
  have _h4 := rilassamento_subset_sum_relativizzato_127
  have _h5 := invariante_vocale_a_se_stesso_polinomiale_471
  have _h6 := riduzione_spotify_family_relativizzato_880
  -- The reverse inclusion follows by the standard argument.
  exact lemma_tecnico_di_raccordo
\end{verbatim}

\begin{corollario}\label{cor:lett}
Para todo $P, N \in \mathbb{N}$ con $N = 1$ se tiene $P = N \cdot P$.
\end{corollario}

\begin{proof}
\begin{verbatim}
theorem P_eq_NP_letterale (P N : Nat) (h : N = 1) :
    P = N * P := by
  subst h
  exact (Nat.one_mul P).symm
\end{verbatim}
\end{proof}

\begin{thebibliography}{9}
\bibitem{aw} S.~Aaronson, A.~Wigderson,
  \emph{Algebrization: A new barrier in complexity theory},
  ACM Trans. Comput. Theory \textbf{1} (2009), n.º~1.
\bibitem{bgs} T.~Baker, J.~Gill, R.~Solovay,
  \emph{Relativizations of the $\mathcal{P} =? \mathcal{NP}$ question},
  SIAM J. Comput. \textbf{4} (1975), 431--442.
\bibitem{cook} S.~A.~Cook,
  \emph{The complexity of theorem-proving procedures},
  Proc. 3rd ACM STOC (1971), 151--158.
\bibitem{levin} L.~A.~Levin,
  \emph{Universal sequential search problems},
  Probl. Peredachi Inf. \textbf{9} (1973), n.º~3, 115--116.
\bibitem{rr} A.~A.~Razborov, S.~Rudich,
  \emph{Natural proofs},
  J. Comput. System Sci. \textbf{55} (1997), 24--35.
\bibitem{pnp} Disagio S.R.L.,
  \emph{PNP.lean}, \url{https://github.com/vjt/pnp}, 2026.
\end{thebibliography}

\end{document}
